\documentclass[hyperref={pdfpagelabels=false},table]{beamer}

\input{lec_common}
\date{}
\begin{document}

\part{Unit 8: More about Plotting}

%\frame{\tableofcontents}


\begin{frame}[fragile]{Object oriented plotting in Python}
\vfill
In this unit we study the different components of a matplotlib plot and
see how to work with them in an object oriented way.
\vfill
We first import the module
\begin{python}
 import matplotlib.pyplot as plt
\end{python}
\end{frame}

\begin{frame}[fragile]{Preparations}
There are two steps:
\begin{itemize}
 \item Construction of the \emph{internal} representation of a plot
 \item Presenting a plot on the screen in a window or in the filesystem as a file in a special graphical file format, e.g.
 pdf, png.
\end{itemize}
The second step requires some preparations, which have to be performed \emph{before} \texttt{pyplot} is imported.
It is a selection of a backend.
\end{frame}
\begin{frame}[fragile]{Backends using Spyder}
 Choose Inline (static) or Automatic (interactive). Then restart the kernel.
 \vfill
\includegraphics[width=.9\textwidth]{inlineOrAutomatic.png}

\end{frame}
\begin{frame}[fragile]{Command driven versus Object oriented Plotting}
\pyth!matplotlib.pyplot! can be used with command-like functions.
\vfill
Various states are preserved.
A state holds information about current figure and current axes.
Under the hood, object oriented programmming (OOP) is used.
\vfill
\pyth!matplotlib! can be also be used for OOP.

Instead of functions, methods of objects are used.
For event driven programs, OOP should be used.
(see next lecture on GUI:s)
\end{frame}
\begin{frame}[fragile]{Example: The command way to plot}
 \begin{python}
x1 = linspace(-2*pi, 2*pi, 100)
x2 = linspace(0, 2, 10)

plot(x1, sin(x1), 'r') # a figure is created
plot(x1, cos(x1), 'b--')
title('My first figure')

figure() # a new figure is created

plot(x2, sqrt(x2), 'go')
plot(x2, x2**2, 'y^')
title('My second figure')
\end{python}
The function calls plot, can be seen as commands. The same command results in different outcomes depending on the state.
\end{frame}
\begin{frame}[fragile]{A second example}
\begin{python}
x1 = linspace(-2*pi, 2*pi, 100)
x2 = linspace(0, 2, 10)

figure() # not needed when using Spyder
subplot(211) # the first subplot using a 2 rows and 1 column grid
plot(x1, sin(x1), 'r')
plot(x1, cos(x1), 'b--')
title('My first subplot')

subplot(212) # the second subplot using a 2 rows and 1 column grid
plot(x2, sqrt(x2), 'go')
plot(x2, x2**2, 'y^')
title('My second subplot')
 \end{python}
 The plots are made in a subplot of the current axes in the current figure.
\end{frame}
\begin{frame}[fragile]{matplotlib as an object-oriented interface}
 For event driven programs you will need references to objects.
 \vfill
\includegraphics[width=.9\textwidth]{figAndAx.png}
\end{frame}
\begin{frame}[fragile]{OOP - getting references to objects}
A single axis object
\begin{python}
 fig, ax = subplots()
\end{python}
\vfill
Multiple axes object
\begin{python}
 fig, ax = subplots(nrows = 2, ncols = 3)
 ax.shape   # (2,3)
 type(ax[0,1])  #  matplotlib.axes._subplots.AxesSubplot
\end{python}

\end{frame}
\begin{frame}[fragile]{Let's compare the two techniques}
Using functions ...
\begin{python}
x = linspace(-2*pi, 2*pi, 100)

figure() # not needed when using Spyder
plot(x, sin(x/2), label = 'sine')
plot(x, cos(x/2), label = '$\cos(x)$')
legend(loc = 'lower center', fontsize = 'small')
\end{python}
Using methods ...
\begin{python}
x = linspace(-2*pi, 2*pi, 100)

fig, ax = subplots()
ax.plot(x, sin(x/2), label = 'sine')
ax.plot(x, cos(x/2), label = '$\cos(x)$')
ax.legend(loc = 'lower center', fontsize = 'small')
\end{python}

\end{frame}
\begin{frame}[fragile]{The figure class}
A figure object is the top level container for all plot elements.
\begin{python}
 fig, ax = subplots()
# plotting code
...
\end{python}
The figure can be saved into a file:
\begin{python}
 fig.figsave("my_figure.png")
\end{python}
Consider also the command \pyth!fig.canvas.get_supported_filetypes()!
\end{frame}
\begin{frame}[fragile]{The axes class}
\includegraphics[width=.7\textwidth]{sphx_glr_anatomy_001.png}\\
\vfill
\hfill {\tiny (Image from \url{https://matplotlib.org/gallery/showcase/anatomy.html})

}

\end{frame}
\begin{frame}[fragile]{Line plots}
\begin{python}
fig, ax = subplots(figsize = (6, 3))

x =linspace(-2*pi, 2*pi, 100)

ax.plot(x, sin(x), 'r--', label = "sine")
ax.plot(x, cos(x), 'b', label = "cosine")

print("The number of line plots is", len(ax.lines))
print("Each line plot is of the type", type(ax.lines[0]))
 \end{python}
 The text output:
\begin{verbatim}
The number of line plots is 2
Each line plot is of the type
         <class 'matplotlib.lines.Line2D'>
\end{verbatim}
\end{frame}
\begin{frame}[fragile]{Getter/Setter Methods}
With the help of getter methods we can get information about the lines (and other objects):
\begin{python}
print("The linestyle of the first line plot is",
                    ax.lines[0].get_linestyle())
print("The color of the last line plot is",
                    ax.lines[-1].get_color())
\end{python}
and we can even extract data from the line
\begin{python}
print("The first y-value of the first plot is ",
                    ax.lines[0].get_ydata()[0])
\end{python}
Setter metods can be used to change attributes, e.g.
\begin{python}
ax.lines[0].set_linewidth(5)
\end{python}
\end{frame}
\begin{frame}[fragile]{Line plots - getting a reference to a line plot}
We can identify a line plot by its label:
\begin{python}
for line in ax.lines:
    if line.get_label() == 'sine':
        line.set_color('green')
\end{python}
\vfill
Note, the plot method returns a list of line plots. In many cases that list has only one element:
\begin{python}
#unpacking a list of one element
li_sin, = ax.plot(x, sin(x))
li_cos, = ax.plot(x, cos(x))
\end{python}
\end{frame}
\begin{frame}[fragile]{Coordinate axes - Labels}
 \begin{python}
fig, ax = subplots()
x =linspace(-2*pi, 2*pi, 100)
ax.plot(x, sin(x))

ax.set_xlim(-3*pi, 3*pi)
ax.set_ylim(-1.2, 1.2)
ax.set_xlabel('x')
ax.set_ylabel('sin(x)')
\end{python}
\end{frame}
\begin{frame}[fragile]{Coordinate axes - Ticklabels}
For publications and slides use few and relevant ticklabels:
\begin{python}
ax.set_xticks([-pi, 0, pi])
ax.set_xticklabels(['-$\pi$', '$0$', '$\pi$'], fontsize = 16)

# set minor ticks
ax.set_xticks([-pi*1.5, -pi*0.5, pi*0.5, pi*1.5], minor = True)

# use empty array to hide ticks
ax.set_yticks([])
\end{python}
\hfill \includegraphics[width=0.4\textwidth]{TickSine.png}
\end{frame}
\begin{frame}[fragile]{Annotations}
Annotations often are more attractive than a legend in the corner
\begin{python}
fig, ax = subplots()
x =linspace(-1.5*pi, 1.5*pi, 100)
ax.plot(x, sin(x))

# we need space above the curve
ax.set_ylim([-1.2, 2])

ax.annotate("local max", xy = (pi/2, 1),
            xytext = (pi, 1.8),
            arrowprops = {'facecolor': 'green',
                          'shrink': 0.1,
                          'width': 2})
\end{python}
\hfill \includegraphics[width=0.3\textwidth]{annotatedSine.png}
\end{frame}
\begin{frame}[fragile]{Filling areas}
Annotations often are more attractive than a legend in the corner
\begin{python}
fig, (ax1, ax2) = subplots(nrows = 1,
                           ncols = 2, figsize = (15, 4))
x =linspace(-2*pi, 2*pi, 100)
y1, y2  = sin(x), cos(x)
ax1.plot(x, y1, 'r')
ax1.plot(x, y2, 'b')
ax1.fill_between(x, y1, y2, color = 'green',
                 alpha = 0.3)
# specify fill condition using where
ax2.plot(x, y1, 'r')
ax2.plot(x, y2, 'b')
ax2.fill_between(x, y1, y2, where = y1 > y2,
                 color = 'red', alpha = 0.3)
\end{python}
\hfill \includegraphics[height=0.222\textheight]{fillSine.png}
\end{frame}
\end{document}
